\section{Introduction}
\label{sec:intro}

% Paragraph 1: Motivation
Does a worker's \textit{first} random access shock matter more than subsequent ones? In tournament-style labor markets, ranking determines access and access determines the opportunity to earn further ranking, so an initial break can trigger a cascade that later opportunities cannot replicate. This paper studies lucky losers in professional tennis: players who lose in the final round of qualifying but enter the main draw when another player withdraws. At Grand Slam tournaments, lucky losers are selected by lottery among the top four ranked qualifying losers, giving a natural experiment in which otherwise similar players are randomly granted or denied a main-draw opportunity. I restrict the analysis to players who have \textit{never previously received} a lucky loser slot, isolating the effect of a first-ever access shock on career trajectories. The setting supports clean measurement of both tournament access and playing ability, a separation that is difficult to obtain in other labor markets.

% Paragraph 2: Mechanism intuition
The mechanism behind career propagation in this setting is path dependence through the ranking system. Professional tennis uses a rolling 52-week ranking that determines direct entry into tournament main draws. A player who earns ranking points at one event may gain direct entry to subsequent events, generating further opportunities to earn points. Access generates exposure, exposure generates accumulation, and accumulation generates future access, forming a feedback loop. A one-time exogenous shock to ranking points can trigger this cascade if the initial shock is large enough. Players who receive a lucky loser slot but lose in the first round earn minimal ranking points and trigger no cascade. Players who win one or more main draw matches earn enough points to cross direct-entry thresholds at lower-tier events, setting off a sequence of additional opportunities. For a player who has never before received this type of opportunity, the first break is potentially pivotal: subsequent opportunities arrive after the player has already crossed direct-entry thresholds, so their marginal value is smaller.

% Paragraph 3: Empirical strategy
The empirical strategy has three components. First, Grand Slam lucky loser allocation is random within the top-four pool, providing the primary identification. I restrict the sample to first-time LL candidates (players with no prior LL awards of any type), so both treated and control players are LL-naive. I estimate immediate effects at the focal event and dynamic outcomes at horizons $h \in \{4, 8, 12, 26, 52\}$ weeks, tracing the trajectory of access and ability from the immediate aftermath through one year. Second, I extend the analysis to non-Grand-Slam events using a control function approach: the generalized residual is computed from the exact enumeration probability that a player's rank among qualifying losers falls within the top slots, a probability determined by the full combinatorial structure of qualifying match pairings. The non-Grand-Slam sample is also restricted to first-time LL candidates. Third, a match-level tournament performance model estimates whether lucky loser entry changes subsequent win probability conditional on opponent strength, again exploiting the Grand Slam lottery for identification.

% Paragraph 4: Main findings
In the first-time ATP Grand Slam sample ($N = 120$; 61 treated, 59 controls), the stacked specification shows that lottery selection raises ranking points by 72.5 at 4 weeks ($p = 0.008$), 82.8 at 12 weeks ($p = 0.002$), and 72.6 at 26 weeks (clustered $p = 0.047$; Fisher $p = 0.164$), fading by 52 weeks. These point estimates are roughly twice the effects estimated on the full sample that pools first-time and repeat LL recipients, consistent with the first break being pivotal. On the WTA tour ($N = 82$), effects concentrate at 8 weeks: $+55.0$ ranking points ($p = 0.017$). The first-time WTA sample also shows Elo gains at 8 weeks ($+15.0$, $p = 0.082$) and 52 weeks ($+36.8$, $p = 0.055$), a pattern absent from the pooled-LL WTA results and consistent with genuine ability improvements on the women's tour. Non-Grand-Slam estimates using the control function approach extend these findings on larger samples ($N = 2{,}054$ ATP; $N = 1{,}423$ WTA): ATP ranking-point gains of 20.6 at 4 weeks ($p = 0.041$) and 44.6 at 26 weeks ($p = 0.002$), together with ATP Elo improvements at 8--26 weeks and a WTA Elo improvement at 12 weeks ($+10.5$, $p = 0.011$). A robustness analysis comparing the marginal effect of the $n$th lucky loser opportunity shows that first-time effects dominate: players receiving their second or third LL slot show weaker and often insignificant gains.

% Paragraph 5: Contribution
This paper makes four contributions. First, I provide the first causal estimates of how a \textit{first} within-career opportunity shock affects trajectories in a tournament labor market. The initial conditions literature has established that graduating in a recession \citep{Oreopoulos2012_graduating_recession}, initial job placement \citep{Oyer2006_initial_placement}, and starting environment \citep{Chetty2014_where_is_land} shape long-run outcomes. I test how a marginal within-career shock, specifically the first such shock, propagates through a transparent tournament market where performance is continuously measured and rankings update weekly. In the settings studied in the initial conditions literature effects persist for a decade or more; here, the access gains fade within a year as the rolling ranking window replaces the initial shock. The roughly twofold increase in point estimates relative to the pooled-LL sample shows that averaging over first-time and repeat recipients attenuates the effect of the pivotal first break; the fadeout indicates that persistence depends on how quickly the market incorporates new performance signals. Second, the access-versus-ability decomposition shows a more nuanced picture than simple ``access begets access.'' While ATP results align with the access channel, WTA first-timers show Elo gains, and non-Grand-Slam first-timers on both tours show ability improvements at medium horizons; a first break can generate learning that subsequent breaks do not replicate. Third, the marginal impact analysis provides direct evidence on diminishing returns to opportunity: Wald tests reject equality of first-LL and second-LL effects for tournament access ($p = 0.002$ ATP; $p = 0.009$ WTA in the non-Grand-Slam sample), and third-plus LL recipients on the WTA tour have negative ranking-point effects at longer horizons. Fourth, the match-level model shows that despite ranking-point gains, first-time ATP GS LL recipients do not outperform controls on a per-match basis: $\hat{\delta} = -0.141$ (player-clustered $p = 0.17$), directionally consistent with overseeding when novice entrants face opponents above their ability level, though the small first-time GS sample cannot separate this from zero. Ranking-point gains therefore come from more draws entered and more matches played rather than a per-match performance improvement. Even so, the access cascade delivers substantial dollar value: the implied total prize-money effect of a first-time Grand Slam LL entry is about \$100{,}000 USD over the subsequent 26 weeks, meaningful against qualifying-level players' annual earnings of \$40{,}000--\$80{,}000.

% Paragraph 6: Roadmap
Section~\ref{sec:literature} reviews the related literature. Section~\ref{sec:data} describes the institutional setting and data. Section~\ref{sec:strategy} presents the econometric framework and identification strategies. Section~\ref{sec:results} reports the main results. Section~\ref{sec:mechanisms} discusses mechanisms. Section~\ref{sec:robustness} presents robustness checks, including the marginal-impact analysis. Section~\ref{sec:conclusion} concludes.
